\scanpage{9}
\mspage{8}
Soit $F$ une figure, \deleted{$\Sigma'$}
\[
\Sigma'_F=\Phi\subset\mathcal{F}_{\text{él}}
\]
la partie de $\mathcal{F}_{\text{él}}$ associée. \deleted{\unsure{S|A}}Alors on a les conditions

a) $\Phi$ est fermée dans l'ens \unsure{ordonné|ordonn} $\mathcal{F}_{\text{él}}$

b) Si \deleted{$\unsure{F_0,F'_0|F_0,F'_2}\ \unsure{\Rightarrow|\ne|\in}$} $F_1,F_2\in\Phi$, \inserted{$F_1\ne F_2$} alors $F_1,F_2$ sont compatibles, i.e. $\Phi$ est une \uline{configuration}

c) $\Phi$ est ``loc. \unsure{finie|fini}'' en un sens convenable.

Donc si on ne veut exclure \unsure{le|les} cas des figures infinies, il faut aussi considérer dans $\mathcal{F}_{\text{el}}$ \unsure{lá|la|là} donnée des parties loc. finies (i.e. dont la famille des \unsure{supports|support} dans $X$ est loc finie) qui forment un antifiltre dans $\mathcal{F}_{\text{el}}$.

Inversement, tout $\Phi$ satisfaisant a) b) c) provient d'une figure et d'une seule.
